%% notes-tex/kg-build-system/sections/4-archive-serving.tex

\section{The archive and the serving hosts\secstatus{todo}}
\label{sec:serving}

\notesec{Taiga is the archive, not a store}
A release directory holds the online backup, the manifest at archive time, the
\texttt{KgRelease} node's properties as \texttt{release.json}, and
\texttt{SHA256SUMS}. \texttt{kg\_release.sh ship} copies it to
\file{/mnt/zhao5/mjvolk3/projects/torchcell/kg-releases/<release>/} on Taiga and
verifies the checksums there; \texttt{list} shows both copies; \texttt{keep} and
\texttt{purge} manage retention. Taiga is where the database lives in the sense
that matters for recovery: every release is a verified set of files that any
host with Neo4j can restore. It is not where a database runs.

\notesec{Why Neo4j does not run a store on NFS}
Neo4j keeps its store as memory-mapped page files and coordinates access with
file locks, and its operations manual lists NFS as an unsupported filesystem for
the data directory\external{Neo4j operations manual, "File system requirements":
NFS is not supported for the database store}. The failure is not theoretical. The
Radiant VM's container serves its old 35-dataset store from the Taiga mount, and
every cold page read raises \texttt{java.io.IOException}; \texttt{ops.sh status}
shows the database online with 7,180,465 nodes and no readable release node. A
restore onto NFS would reproduce that fault for the new store.

\notesec{The deploy step}
\texttt{kg\_release.sh deploy <release> [--pin]} runs on the host that will serve.
It pulls the release directory from Taiga when the local archive root lacks it,
verifies \texttt{SHA256SUMS}, restores the artifact with \texttt{neo4j-admin
database restore} into a new database \texttt{kg-<release>} while the DBMS stays
online, brings the database up with \texttt{CREATE DATABASE ... WAIT}, compares
the restored node's release id, version and dataset count with
\texttt{release.json}, and only then retargets the \texttt{latest} alias (and
\texttt{pinned} with \texttt{--pin}). The previous database stays for rollback.
It ends with \texttt{ops.sh sync}, which exits 1 while two hosts serve different
releases or any host serves a store without a release node.
\texttt{restore-test <release>} is the same restore and the same checks into a
throwaway database, dropped at the end; it is what proved the 3.0 artifact
restorable (10 min 13 s, counts identical to the served store).

\notesec{The fix for DIVERGED}
Radiant has 16 CPUs, 62 GB of memory, a 40 GB root disk and the Taiga mount, and
no other disk. Until it has a local block volume for the container's
\texttt{/data} there is no Neo4j-side fix, and \texttt{ops.sh sync} stays
DIVERGED by design. The fix is two steps.

\begin{enumerate}
\item Attach a block volume and mount it at the path the container maps to
  \texttt{/data}. The restored 3.0 store is 151 GB, so 500 GB holds two releases
  side by side with headroom; the earlier 2 TB estimate predates the store
  shrinking to this size.
\item On Radiant, with the archive root pointed at Taiga:
\begin{verbatim}
ARCHIVE_ROOT=/mnt/zhao5/mjvolk3/projects/torchcell/kg-releases \
NEO4J_CONTAINER=tc-neo4j \
  bash scripts/kg_release.sh deploy 2026.10.06-4b293d34 --pin
\end{verbatim}
\end{enumerate}

After that, every release is \texttt{ship} then \texttt{deploy}, and the served
version on Radiant follows the archive. The one interim alternative is a bolt
proxy on Radiant forwarding port 7687 to GilaHyper, which gives clients one
hostname but keeps the live store on GilaHyper rather than Taiga; it has not been
set up.
